\chapter*{ВВЕДЕНИЕ}
\addcontentsline{toc}{chapter}{ВВЕДЕНИЕ}

\section*{Актуальность}
Существует задача автоматического подбора рецепта, соответствующего предпочтениям пользователя, по заданному набору ингредиентов.
Её решение актуально в случае, когда человеку разрешается питаться ограниченным набором продуктов по медицинским причинам.
Разработанная программа предоставляет пользователю возможность ввести список доступных и разрешённых ему ингредиентов и получить список рецептов, которые содержат только их.
Таким образом, человек находит новые для себя рецепты, которые соответствуют его диете и при этом нравятся ему.
Следовательно, пользователю становится легче придерживаться определённой диеты, что позитивно сказывается на его здоровье и самочувствии~\cite{zholobova2022nutrition}.
Более того, рекомендательная система рецептов блюд позволяет экономить время на выбор рецепта, так как человеку не нежно самостоятельно искать рецепты и анализировать их на соответствие своим потребностям.
Также приложение предоставляет подробные инструкции приготовления блюд.
Это даёт пользователю возможность приготовления блюд без соответствующих знаний и навыков.
Исходя из представленных рассуждений, рассматриваемая задача является актуальной.

\section*{Цель работы}
\textbf{Целью} данной работы является разработка системы рекомендаций рецептов блюд, учитывающей как предпочтения пользователя, так и его возможности.
Под возможностями понимается наличие необходимых ингредиентов.

Для достижения поставленной цели необходимо выполнить следующие \textbf{задачи}:
\begin{itemize}
    \item проанализировать предметную область алгоритмов рекомендаций;
    \item описать основные методы формирования рекомендаций и выполнить их классификацию;
    \item проанализировать существующие программы для рекомендаций рецептов;
    \item выбрать метод рекомендаций рецептов;
    \item спроектировать рекомендательную систему в виде программного обеспечения;
    \item реализовать программное обеспечение;
    \item реализовать выбранные алгоритмы рекомендаций;
    \item подготовить набор данных для обучения и тестирования моделей рекомендаций;
    \item сравнить точность реализованных моделей рекомендаций на подготовленном наборе данных.
\end{itemize}
